\documentclass[../../main.tex]{subfiles}
\graphicspath{{../../image/}{../../figure/}}
\begin{document}

\section{注记}\label{section:ch12-注记}

关于真非线性两守恒律组理论的各个方面，已有大量文献。本章通过广义特征线理论处理问题的途径由作者首创。这些结果中的许多（偶尔在稍强假设下）更早是在由随机选取方法构造的解的框架下导出的，该方法将在第 XIII 章中介绍。
该方向的开创性贡献是 Glimm 与 Lax [1]，他们引入并使用了"近似特征线"与"近似守恒律"的概念。近似特征线，即经典特征线与激波的拼接，与本书意义下的广义特征线密切相关。然而两种途径之间有根本区别：Glimm 与 Lax 沿近似特征线前向传播跟踪，而此处我们回顾性地看待广义特征线。

\cref{section:ch12-熵熵流对与速矢变换}中讨论的 Lax 熵最早由 Lax [4] 引入。Hodograph 变换由 Riemann [1] 发现。关于详细讨论及其在空气动力学中的应用，见 Courant 与 Friedrichs [1] 及 von Mises [1]。关于在数学物理其他领域的应用，见 D. Fusco [1]。

\cref{theorem:ch12-局部正则存在}的一个稍强版本由 DiPerna [3] 对随机选取方法构造的解建立。定理 12.4.1 改进了 Dafermos [19] 中的一个命题。

定理 12.5.1、12.6.1 与 12.6.4 最初由 Glimm 与 Lax [1] 对随机选取方法构造的解建立。
此处的处理采用并精化了 Dafermos 与 Geng [1,2] 对特殊系统、Trivisa [1] 对一般系统（尽管解"可数正则"时）所发展的方法。Trivisa [2] 将这些结果推广到配备 Riemann 不变量坐标系的真非线性 $n$ 守恒律组。

Tsikkou [1,2] 用同一方法处理等熵弹性的特殊系统 (4.1.11)。
她用分段线性函数逼近 $\sigma(u)$，对所得特征线模式做详尽考察，从而导出总变差的尖锐界。她还表明，当 $\sigma'(u)$ 含跳跃间断时，解的总变差可能在有限时间内爆破。

\cref{section:ch12-解的正则性}的结果更早由 DiPerna [3] 对随机选取方法构造的解建立。

对初值在 $L^1$ 中的解，Temple [5] 导出 $O(1/\sqrt{\log t})$ 的衰减速率。定理 12.8.1 中建立的 $O(t^{-\frac12})$ 衰减速率取自 Dafermos [19]，是尖锐的。
类似地，引理 12.8.2 改进了 Temple [8] 的较早结果。$L^1$ 稳定性现已对一般系统建立；见第 XIV 章。

生成 $N$ 波剖面的机制很早就通过形式渐近被理解（见 Courant 与 Friedrichs [1]），尽管缺乏严格证明（Lax [2]）。在 DiPerna [4,6] 与 Tai-Ping Liu [8,9,22] 的一系列论文中，对由随机选取方法构造的具紧支集初值的解，以逐渐尖锐的速率建立了到 $N$ 波的衰减，不仅对真非线性两守恒律组，甚至对特征族为真非线性或线性退化的 $n$ 守恒律组。关于通过变量替换解耦特征场的另一途径，见 Ayad 与 Heibig [1,2]。定理 12.9.1 中记录的衰减速率是尖锐的。
当初值不具紧支集、而是当 $x\to\pm\infty$ 时逼近不同极限 $U_L$ 与 $U_R$ 时，解 $U$ 当 $t\to\infty$ 收敛到具初值 (9.1.12) 的 Riemann 问题解；见 Tai-Ping Liu [6] 并与\cref{section:ch12-L1估计}中讨论的标量情形比较。
对非真非线性系统所知相对甚少；见 Zumbrun [1,2]。

定理 12.10.1 归功于 Glimm 与 Lax [1]，而定理 12.10.2 取自 Dafermos [21]。
关于相对论等熵气流系统周期解的稳定性，见 Calvo、Colombo 与 Frid [1]。具周期初值的解的衰减可能是两守恒律组特有的。事实上，R. Young [3,4] 与 Temple 及 Young [3,4] 的工作表明，对非等熵气流系统，具周期初值的解保持有界但不一定衰减。另一方面，Qu 与 Xin [1] 的近期工作表明此类解的寿命非常长。另一方面，是或可变换为线性退化的系统通常容许光滑时间周期解。例子见 M. Shearer [7]。

关于特征线理论在研究具重合激波与稀疏波曲线的特殊系统（Temple [3]）解的唯一性、正则性与大时间行为中的应用，见 Serre [7,11]、Dafermos 与 Geng [1,2]、Heibig [2]、Heibig 与 Sahel [1] 及 Ostrov [1]。此类系统的 $BV$ 解已由 Godunov 差分格式（LeVeque 与 Temple [1]）以及消黏性方法（Serre [1,11]）构造。关于初边值问题解在 $L^\infty$ 中的爆破，见 Bourdarias、Gisclon 与 Junca [2]。

关于通过 Glimm-Lax 途径处理两守恒律组的文献十分丰富。关于将 Lax-Glimm 理论在激波与稀疏波曲线较少限制下推广到真非线性系统，见 Bianchini、Colombo 与 Monti [1,2]。关于非齐次流系统的例子，见 Frid、Risebro 与 Sande [1]。

\end{document}
